% Generated draft: requires mathematical and source review.
\documentclass[11pt]{article}
\usepackage[margin=25mm]{geometry}
\usepackage{fontspec}
\setmainfont{Latin Modern Roman}
\usepackage{amsmath,amssymb,mathtools,unicode-math}
\setmathfont{Latin Modern Math}
\usepackage{enumitem,needspace,xurl,hyperref}
\setlength{\parindent}{0pt}
\setlength{\parskip}{5pt}
\setlength{\emergencystretch}{3em}
\newcommand{\N}{\mathbb{N}}
\newcommand{\Z}{\mathbb{Z}}
\newcommand{\Q}{\mathbb{Q}}
\newcommand{\R}{\mathbb{R}}
\newcommand{\C}{\mathbb{C}}
\newcommand{\F}{\mathbb{F}}
\newcommand{\A}{\mathbb{A}}
\newcommand{\PP}{\mathbb{P}}
\newcommand{\E}{\mathbb{E}}
\newcommand{\eps}{\varepsilon}
\newcommand{\dd}{\,\mathrm{d}}
\begin{document}
\section*{UnsolvedMath ID: 30002592}
\textbf{Source problem code:} catalogue entry
\subsection{Definitions and mathematical statement}
For an integer \(n\geq 1\), write
\[
\operatorname{Conf}_n(\C)=\{(z_1,\ldots,z_n)\in\C^n:z_i\neq z_j\text{ whenever }i\neq j\}.
\]
Using the usual braid-group conventions, let \(B_n=\pi_1(\operatorname{Conf}_n(\C)/S_n)\) be the Artin braid group and let
\[
PB_n=\pi_1(\operatorname{Conf}_n(\C))\cong\ker(B_n\longrightarrow S_n)
\]
be the pure braid group, with compatible basepoints understood. These definitions fix notation; the supplied record does not specify the intended range of \(n\) for its questions.

The supplied record asks whether a deformation \(K\)-theory identification can be extended to pure braid groups. Its displayed target, after substituting \(G=PB_n\), is the schematic expression
\[
K_{\mathrm{def}}(PB_n)\stackrel{?}{\cong}H_*(\operatorname{Irr}(PB_n)).
\]
Here \(K_{\mathrm{def}}\) denotes deformation \(K\)-theory in the unspecified convention of the cited source, and \(\operatorname{Irr}(G)\) is described by the record as the topological monoid of irreducible representations of \(G\).

\textbf{Unresolved formulation qualifications.} The record does not specify the representations' coefficient field, whether they are unitary, their allowed dimensions, the equivalence relation on representations, or the topology on the resulting representation space. Moreover, irreducible representations alone are generally not closed under direct sum. Thus the phrase ``topological monoid of irreducible representations'' does not determine a monoid: it could require a construction generated by irreducibles or another operation, neither of which is supplied. The homology coefficients, grading conventions, and meaning of the displayed isomorphism are also absent. In particular, the notation does not distinguish a comparison involving a spectrum from one involving its homotopy groups. The theorem to be extended, its original hypotheses, and its comparison map are not stated. Consequently, this display records the requested extension but does not yet define a precise isomorphism conjecture.

The original statement additionally asks whether the corresponding deformation \(K\)-theory and complex topological \(K\)-theory for the full Artin braid groups \(B_n\) can be identified. The supplied background also mentions an associated representation-ring construction for these groups. Neither the space or algebra whose complex topological \(K\)-theory is intended, nor the representation ring, grading, comparison map, or proposed identification is specified. This additional question must therefore remain unresolved at the formulation level; it cannot be replaced by a particular comparison theorem from the supplied information.
\subsection{Short English statement}
The problem asks whether an unspecified comparison between deformation \(K\)-theory and homology built from irreducible representations extends to pure braid groups. The original version also asks for a related identification for full Artin braid groups. Essential definitions and comparison data are missing, so both questions require clarification before they become precise mathematical conjectures. The supplied open status is a dated curation label.
\subsection{Sources}
\begin{itemize}
\item Local registry: \texttt{data/unsolvedmath/problems.json}, with supplementary registry \texttt{catalog/problems.json}; supplied complete record for permanent UnsolvedMath ID \(30002592\), titled ``Deformation K-Theory of Pure Braid Groups,'' exact source problem code \texttt{catalogue entry}. The record gives proposed year 2014 and names no proposer.
\item Supplied report citation: \emph{Algebraic Structures in Low-Dimensional Topology} (2015), Oberwolfach Reports, pp.~1403--1458, DOI \texttt{10.4171/owr/2014/26}. \url{https://doi.org/10.4171/owr/2014/26}. The record describes Oberwolfach Reports as documenting workshops and their open-problem discussions. The report citation is supplied provenance; its contents have not been independently consulted for this draft.
\item Supplied literature-source attribution: Bryden, \emph{Algebraic Structures in Low-Dimensional Topology}. \url{https://doi.org/10.4171/owr/2014/26}. The supplied assessment attributes excision and a free-group formula for pure-braid deformation \(K\)-theory to Bryden; the exact results and hypotheses are not included in the record.
\item Supplied literature source: Azzali et al., \emph{K-theory of pure braid groups}. \url{https://arxiv.org/abs/2105.07848}.
\item Literature assessment dated exactly \texttt{2026-08-21}: the supplied curation labels the problem open and reports that the extracted extension problem and associated full Artin-braid representation-ring construction remained unresolved in its literature search. This is supplied dated curation, not an independently verified literature assessment. The record expressly requests independent verification of time-sensitive status and cited claims.
\item The supplied statement status is \texttt{corrected\_verified}, accompanied by a claim that the clean statement was checked against the official report and restores its hypotheses and notation. This is a supplied verification claim, not verification performed for this draft.
\end{itemize}
\end{document}
